% ============================================================================
% SEC04.TEX - Section 7.4: Parameters
% ============================================================================

\section{Parameters}

\begin{learningoutcomes}
\item Menggunakan Animator parameters
\item Mengset parameters melalui script
\item Menggunakan trigger untuk one-time events
\end{learningoutcomes}

\subsection{Jenis Parameters}

\begin{itemize}
\item \textbf{Float}: Nilai desimal (speed, health percentage)
\item \textbf{Int}: Nilai bulat (state index, count)
\item \textbf{Bool}: True/false (isGrounded, isAttacking)
\item \textbf{Trigger}: Event sekali pakai (jump, attack, die)
\end{itemize}

\subsection{Menggunakan Parameters dari Script}

\begin{lstlisting}[language={[Sharp]C}]
using UnityEngine;

public class AnimatorController : MonoBehaviour
{
    private Animator animator;
    
    void Start()
    {
        animator = GetComponent<Animator>();
    }
    
    void Update()
    {
        // Set float parameter
        float speed = GetSpeed();
        animator.SetFloat("Speed", speed);
        
        // Set bool parameter
        bool isGrounded = IsGrounded();
        animator.SetBool("IsGrounded", isGrounded);
        
        // Set trigger
        if (Input.GetKeyDown(KeyCode.Space))
        {
            animator.SetTrigger("Jump");
        }
    }
}
\end{lstlisting}

\begin{catatan}
Parameters menghubungkan gameplay dengan animasi. Update parameters setiap frame di Update() untuk responsif. Gunakan trigger untuk events sekali pakai seperti jump atau attack.
\end{catatan}
